% !TEX root = ../projekt.tex
% =============================================================================
% STAV: pracovný podklad. Prepísať vlastnými slovami, doplniť citácie.
% =============================================================================
\chapter{Úvod}
\label{chap:uvod}

Rýchlosť počítačových sietí rastie rýchlejšie než výkon jedného jadra procesora.
Pri rýchlostiach 100\,Gb/s a viac už jedno jadro nestíha spracovať všetky prijaté
pakety, preto sieťová karta rozdeľuje prevádzku do viacerých DMA kanálov a každý
kanál obsluhuje iné jadro. \TODO{citácia a konkrétne čísla: čas na jeden paket pri 100\,Gb/s}

Štandardným mechanizmom tohto rozdelenia je \emph{Receive Side Scaling} (RSS).
Z~hlavičky paketu vypočíta hash a podľa neho vyberie DMA kanál. Všetky pakety jedného
toku tak skončia v~tom istom kanáli, čo zachová ich poradie. Cenou za túto vlastnosť je
nerovnomerná záťaž: toky majú veľmi rozdielnu veľkosť a kanál, na ktorý padne niekoľko
veľkých tokov, je preťažený, zatiaľ čo iné kanály majú voľnú kapacitu.

Cieľom práce je navrhnúť a analyzovať varianty distribúcie paketov do DMA kanálov
jednotkou RSS frameworku NDK, vyvíjaného združením CESNET. Práca uvažuje úpravy, ktoré
zachovávajú pravidlo ,,jeden tok = jeden kanál``, aj úpravy, ktoré rozdeľujú pakety
jedného toku na podmnožinu kanálov. Na ich vyhodnotenie vznikol simulátor jednotky RSS
pre najviac 128 DMA kanálov, ktorý pracuje so vzorkami reálnej sieťovej prevádzky.

\TODO{Sem patrí 2--3 vety o hlavnom výsledku, až keď budú hotové varianty z~kapitoly~\ref{chap:vyhodnotenie}.}

\paragraph{Štruktúra práce.}
Kapitola~\ref{chap:ndk} predstavuje platformu NDK. Kapitola~\ref{chap:rss} popisuje
mechanizmus RSS. Kapitola~\ref{chap:metodika} zavádza metriky a metodiku merania
a kapitola~\ref{chap:simulator} implementovaný simulátor. Kapitola~\ref{chap:analyza}
analyzuje východiskové RSS a ukazuje, čo jeho nerovnomernosť spôsobuje.
Kapitola~\ref{chap:varianty} navrhuje varianty úprav, kapitola~\ref{chap:vyhodnotenie}
ich vyhodnocuje a kapitola~\ref{chap:zaver} prácu uzatvára.
